%==========================================================
% B201 Computer Science Lab - Individual Project Report
% Eser Menke
% Written in LaTeX (compiles with pdfLaTeX)
%==========================================================
\documentclass[11pt,a4paper]{article}

\usepackage[utf8]{inputenc}
\usepackage[T1]{fontenc}
\usepackage[margin=2.5cm]{geometry}
\usepackage[dvipsnames]{xcolor}
\usepackage{titlesec}
\usepackage{enumitem}
\usepackage{parskip}
\usepackage{hyperref}

\definecolor{accent}{HTML}{2B6CB0}

\hypersetup{
  colorlinks=true,
  urlcolor=accent,
  linkcolor=accent,
  citecolor=accent
}

\titleformat{\section}{\large\bfseries\color{accent}}{\thesection.}{0.5em}{}
\titlespacing{\section}{0pt}{14pt}{6pt}
\titleformat{\subsection}{\normalsize\bfseries}{}{0em}{}
\titlespacing{\subsection}{0pt}{10pt}{4pt}

\begin{document}

%========================== TITLE ===========================
\begin{center}
  {\large Gisma University of Applied Sciences}\\[2pt]
  {\normalsize School of Computer Science}\\[14pt]
  {\Large\bfseries\color{accent} B201 Computer Science Lab}\\[6pt]
  {\Large\bfseries Computer Science Portfolio Report}\\[14pt]
  {\normalsize Eser Menke}\\[2pt]
  {\small Individual Project}\\[2pt]
  {\small June 2026}
\end{center}

\vspace{10pt}
\hrule
\vspace{14pt}

%======================= INTRODUCTION =======================
\section{Introduction}

This report is about the computer science portfolio I built for the B201 Computer Science
Lab module. The task was to create a portfolio that shows my technical skills, as if I were
applying for a working student position. So I made a personal website that introduces me,
lists the skills I am starting to learn, and shows a few small projects I have worked on.

I am in my first semester, so I did not try to make the portfolio look like the work of an
experienced developer. Instead I tried to keep it simple, clear, and honest about where I am
right now. The two main links for this project are below.

\begin{itemize}[leftmargin=1.4em, itemsep=2pt, topsep=2pt]
  \item \textbf{Portfolio website:} \url{https://esermenke.github.io}
  \item \textbf{GitHub repository:} \url{https://github.com/EserMenke/EserMenke.github.io}
\end{itemize}

%==================== WEBSITE STRUCTURE =====================
\section{Structure of the Website}

The website is a single page that you scroll through, with a menu at the top that jumps to
each part. I chose a single page instead of many separate pages because the portfolio does
not have a lot of content yet, and one clear page felt easier to read and easier to look
after. The page is split into these sections:

\begin{itemize}[leftmargin=1.4em, itemsep=2pt, topsep=2pt]
  \item \textbf{Header / navigation:} my name and a small menu (About, Skills, Projects, CV, Contact).
  \item \textbf{Hero:} a short line that says who I am, with two buttons.
  \item \textbf{About:} a short paragraph about me and why I chose software engineering.
  \item \textbf{Skills:} the languages and tools I have started using, each marked as beginner.
  \item \textbf{Projects:} the small things I have built so far.
  \item \textbf{CV:} a button to download my CV as a PDF.
  \item \textbf{Contact:} my email, GitHub, and LinkedIn.
\end{itemize}

%================== REPOSITORY STRUCTURE ====================
\section{Structure of the GitHub Repository}

The code lives in a public GitHub repository. I named the repository
\texttt{EserMenke.github.io} on purpose, because GitHub Pages uses that name automatically
to publish the site at \url{https://esermenke.github.io}. The repository is
organised like this:

\begin{itemize}[leftmargin=1.4em, itemsep=2pt, topsep=2pt]
  \item \texttt{index.html} \quad the main page with all the sections.
  \item \texttt{style.css} \quad all of the styling for the page.
  \item \texttt{cv/} \quad my CV (the PDF and the LaTeX source).
  \item \texttt{projects/} \quad source files for the exercises shown on the site.
  \item \texttt{assets/} \quad images and other static files.
  \item \texttt{README.md} \quad a short explanation of the project.
\end{itemize}

I kept the styling in its own file instead of writing it inside the HTML, because it keeps
the page easier to read and easier to change later.

%==================== TOOLS AND TECH ========================
\section{Tools and Technologies}

I built the site with plain \textbf{HTML} and \textbf{CSS}, and I did not use any JavaScript.
I made that choice because I am still a beginner, and I wanted to fully understand everything
on the page rather than copy in code I could not explain. For the structure and the styling I
used the MDN Web Docs and W3Schools as references when I got stuck (Mozilla, 2024; W3Schools,
2024). For the text I used the Inter font from Google Fonts (Google Fonts, 2024).

For version control and hosting I used \textbf{Git}, \textbf{GitHub}, and \textbf{GitHub
Pages}, which let me put the site online for free straight from the repository. My
\textbf{CV} and this \textbf{report} were both written in \textbf{LaTeX}, which I compiled
in Overleaf (Overleaf, 2024).

%==================== DESIGN DECISIONS ======================
\section{Design Decisions}

Since I am still learning, most of my design choices came from small tricks I picked up here
and there from tutorials and examples, and I tried to use them in a way that made sense for my
own page. A few of the main choices were:

\begin{itemize}[leftmargin=1.4em, itemsep=2pt, topsep=2pt]
  \item \textbf{One accent colour.} I used a single blue for links, buttons, and section
        titles, so the page looks calm and consistent instead of busy.
  \item \textbf{Cards for skills and projects.} Putting each skill and project in its own
        box makes the information easier to scan.
  \item \textbf{Responsive layout with CSS only.} The cards and the menu change to fit smaller
        screens, so the site still works on a phone.
  \item \textbf{Honest content.} I marked every skill as beginner and kept the projects small,
        because I wanted the portfolio to match where I actually am.
\end{itemize}

I also tried to keep the writing on the site short and plain, so it is easy to read.

%================= STRENGTHS AND WEAKNESSES =================
\section{Strengths and Weaknesses}

\subsection{Strengths}
The site is simple and clean, it loads fast, and it works on both computers and phones.
Because it only uses HTML and CSS, I can explain every part of it, and it is easy to update
when I learn something new or finish a new project.

\subsection{Weaknesses}
The site is also quite basic. Without any JavaScript it has no interactive features, the
design is plain, and the projects are small beginner exercises. These are honest limits of
being in my first semester, and they are the parts I most want to improve over time.

%======================= REFLECTION =========================
\section{Reflection}

The part I both enjoyed the most and struggled with the most was the C++ project that I made
as an example for the portfolio. I really tried to build it well and keep improving it. Even
installing C++ and getting it set up was a challenge on its own, but the whole journey turned
out to be pretty fun.

While working on this assignment I also learned several new things. I found out what LaTeX
actually is and how to use it for my CV and this report, I got my first real experience with
C++, and I picked up a lot of design ideas while putting the project together.

%==================== FUTURE IMPROVEMENTS ===================
\section{Future Improvements}

There are a few things I would like to add as I get better. Most of all, I would really love
to add a smooth, fluid interface to the site, but I have not figured out how to do that yet.
Beyond that, I want to add more and bigger projects over time, include real screenshots and
short write-ups for each one, and slowly improve the overall design as my skills grow.

%======================= CONCLUSION =========================
\section{Conclusion}

This project gave me my first full experience of building something from start to finish:
designing a website, writing it in HTML and CSS, putting it online with GitHub Pages, and
documenting it with LaTeX. The result is simple, but it honestly shows where I am as a
first-year student, and it gives me a base I can keep building on.

%======================= REFERENCES =========================
\section*{References and Bibliography}
\addcontentsline{toc}{section}{References and Bibliography}
\begin{itemize}[leftmargin=1.4em, itemsep=4pt, topsep=4pt, label={}]
  \item GitHub (2024) \textit{GitHub Pages documentation}. Available at:
        \url{https://docs.github.com/en/pages} (Accessed: 26 June 2026).
  \item Google Fonts (2024) \textit{Inter}. Available at:
        \url{https://fonts.google.com/specimen/Inter} (Accessed: 26 June 2026).
  \item Mozilla (2024) \textit{HTML: HyperText Markup Language}. MDN Web Docs. Available at:
        \url{https://developer.mozilla.org/en-US/docs/Web/HTML} (Accessed: 26 June 2026).
  \item Overleaf (2024) \textit{Learn LaTeX in 30 minutes}. Available at:
        \url{https://www.overleaf.com/learn/latex/Learn_LaTeX_in_30_minutes} (Accessed: 26 June 2026).
  \item W3Schools (2024) \textit{HTML Tutorial}. Available at:
        \url{https://www.w3schools.com/html/} (Accessed: 26 June 2026).
\end{itemize}

\end{document}
